% ====================================================================
% THESIS SECTION: GENETIC ALGORITHM EXPERIMENTAL RESULTS AND EVALUATION
% This file has been automatically generated and populated with real metrics.
% ====================================================================

\section{Genetic Algorithm Experimental Results and Evaluation}
\label{sec:ga_experimental_results}

In this section, the performance of the developed Pure Genetic Algorithm (GA) with Heuristic Initialization for the Exam Timetabling Problem (ETP) is evaluated. Using the large-scale international competition benchmark dataset \textbf{agh-fal17}, we analyze the search performance, constraint resolution capability, and convergence behavior.

\subsection{Experimental Configuration Setup}
\label{subsec:experimental_configurations}

To optimize the exam schedule under a highly constrained search space, we configured the GA with High Selection Pressure and Heuristic Initialization:
\begin{itemize}
    \item \textbf{Population Size:} $100$
    \item \textbf{Mutation Rate:} $0.10$
    \item \textbf{Tournament Selection Size:} $15$
    \item \textbf{Initialization Scheme:} Stochastic Heuristic Initialization (ordering classes by student conflict degree and selecting candidates to minimize initial conflicts)
\end{itemize}

The algorithm was executed for $5$ generations.

\subsection{Performance Analysis and Quantitative Results}
\label{subsec:performance_comparisons}

The numerical results obtained from the run are summarized in Table~\ref{tab:ga_experiments}. The table presents the execution runtime (seconds), the best final fitness value, and the corresponding student, room, time, and distribution penalty breakdowns.

\begin{table}[htbp]
\centering
\caption{Performance Analysis of the Heuristically Initialized Genetic Algorithm}
\label{tab:ga_experiments}
\begin{tabular}{lr}
\hline
\textbf{Metric} & \textbf{Value} \\ \hline
Population Size & 100 \\
Mutation Rate & 0.10 \\
Tournament Size & 15 \\
Execution Time (s) & 586.06 \\
Student Penalty & 43663 \\
Room Penalty & 145304 \\
Time Penalty & 1493 \\
Distribution Penalty & 2657 \\ \hline
\textbf{Total Penalty} & \textbf{409446} \\ \hline
\end{tabular}
\end{table}

The heuristic initialization allowed the GA to start with a highly optimized initial population, drastically bypassing the chaotic early generations of random search. This allowed the high selection pressure to focus on fine-tuning soft constraints.

\subsection{Convergence Curve Analysis}
\label{subsec:convergence_curve_analysis}

The evolutionary trajectory and optimization trends over generations are illustrated in the convergence chart in Figure~\ref{fig:ga_convergence_curves}.

\begin{figure}[htbp]
\centering
\includegraphics[width=0.85\textwidth]{experiment_results.png}
\caption{Generation-by-Generation GA Convergence Curve}
\label{fig:ga_convergence_curves}
\end{figure}

The convergence profile indicates a smooth descent, demonstrating the stabilizing effect of starting with a heuristically guided population. The algorithm managed to balance room allocation, timeslot availability, and student schedule spreads to achieve a low-conflict timetable of \textbf{409446}.

